%% en/sections/02_regole.tex — Rules (full EN)

\chapter{The Rules}
\markboth{The Rules}{The Rules}

\lettrine[lines=3]{V}{espri} 1282 uses the \textbf{Forged in the Dark}
(FitD) system, created by John Harper for \textit{Blades in the Dark}.
This section presents the rules adapted to the historical context and
multi-table structure. Experienced FitD players will find the system
familiar; new players will find everything they need here.

\section{The Heart of the System}

When a character attempts something risky and the outcome is uncertain,
roll dice. Gather a number of six-sided dice equal to the relevant action
rating, roll, and consider the highest result.

\regola[Action Roll]{
  \textbf{6}: Full success. You get what you want, without cost.\\
  \textbf{4--5}: Partial success. You get what you want, but with
  a complication, a cost, or reduced effect.\\
  \textbf{1--3}: Failure. Things go wrong. The GM introduces a
  negative consequence.\\[4pt]
  If you have no dice (action at 0), roll \textbf{two dice} and take
  the lower result.
}

\subsection{Position and Effect}

Before each roll, the GM establishes two fundamental parameters:

\textbf{Position} --- your exposure to risk:
\begin{itemize}[nosep]
  \item \textit{Controlled} --- you are safe, risk is minimal
  \item \textit{Risky} --- tense situation, consequences are serious
  \item \textit{Desperate} --- you are in real danger, consequences
    can be devastating
\end{itemize}

\textbf{Effect} --- how much you can impact the situation:
\begin{itemize}[nosep]
  \item \textit{Limited} --- your impact is small
  \item \textit{Standard} --- your impact is as expected
  \item \textit{Great} --- your impact exceeds expectations
\end{itemize}

\nota[GM Note]{
  Position and effect are narrative tools, not fixed parameters.
  Change them freely in response to player actions. A well-elaborated
  plan can shift position from Risky to Controlled; a sudden distraction
  can do the reverse.
}

\secsep

\section{The Actions}

Characters in \textit{Vespri 1282} act through \textbf{twelve actions}
grouped into four attributes reflecting the tensions of the historical period.

\subsection{Heart --- \textit{Passion and Will}}

\begin{description}[leftmargin=2em, style=nextline]
  \item[\cinzel Incite] Inflame minds, move crowds, make emotional appeals.
    The action of the preacher, the demagogue, the martyr.
  \item[\cinzel Resist] Hold firm under pressure, endure physical or emotional
    pain, refuse to break under interrogation.
  \item[\cinzel Defy] Confront an adversary directly, put yourself before
    danger, accept open confrontation.
\end{description}

\subsection{Mind --- \textit{Intelligence and Cunning}}

\begin{description}[leftmargin=2em, style=nextline]
  \item[\cinzel Study] Gather information, read people and situations,
    understand others' plans.
  \item[\cinzel Plot] Elaborate plans, build networks, coordinate agents
    at scale.
  \item[\cinzel Deceive] Lie convincingly, construct false identities,
    confuse adversaries.
\end{description}

\subsection{Body --- \textit{Strength and Dexterity}}

\begin{description}[leftmargin=2em, style=nextline]
  \item[\cinzel Act] Move quickly, perform precise physical actions,
    move unseen.
  \item[\cinzel Fight] Use violence effectively, protect allies, disarm
    or neutralise opponents.
  \item[\cinzel Manoeuvre] Move through controlled spaces, pass unnoticed,
    escape surveillance.
\end{description}

\subsection{Spirit --- \textit{Influence and Faith}}

\begin{description}[leftmargin=2em, style=nextline]
  \item[\cinzel Command] Lead with authority, give orders that are followed,
    maintain discipline.
  \item[\cinzel Persuade] Convince with arguments, negotiate, seduce with words.
  \item[\cinzel Invoke] Appeal to greater forces --- God, justice, the memory
    of ancestors. Special action of the Clergy, available to all in extreme moments.
\end{description}

\secsep

\section{Stress and Trauma}

Each character has \textbf{9 Stress boxes}. Stress accumulates when
you suffer consequences, use Fortune, or put yourself in desperate
situations. When the ninth box fills, the character suffers a
\textbf{Trauma} and the track clears.

\subsection{Traumas of Vespri 1282}

\begin{itemize}[nosep]
  \item \textit{Distrust} --- you no longer trust anyone, not even allies
  \item \textit{Obsession} --- one goal consumes every other consideration
  \item \textit{Guilt} --- you did something unforgivable; you cannot stop
    thinking about it
  \item \textit{Fear} --- a specific type of situation paralyses you
  \item \textit{Lost Faith} --- you no longer believe in what brought you here
    (God, the cause, your lord)
\end{itemize}

A character with \textbf{4 traumas} is out of play: they withdraw,
collapse, or betray.

\regola[Using Fortune]{
  Before a roll, declare you are using Fortune. Add \textbf{1d6} to
  the roll. Then suffer \textbf{2 Stress}. You may use Fortune up to
  three times on the same roll.
}

\regola[Resisting Consequences]{
  Roll dice for the relevant attribute. Each result above 3 reduces
  the severity of the consequence. A 6 eliminates it entirely. You
  still suffer \textbf{1 Stress} for resisting, regardless of result.
}

\secsep

\section{Multi-Table Mechanics}

\subsection{Faction Score}

Each faction accumulates a score from 0 to 4 during the session.
The total score across all four tables (0--16) determines the final
outcome of the Vespers night.

\nota[For the Coordinating GM]{
  The coordinating GM tracks all table scores on a central sheet.
  Scores are not revealed to players during the session: they measure
  what is happening off-screen. Each table only knows what its own
  GM communicates narratively.
}

\subsection{Conjunction Moments}

At three moments during the session the coordinating GM distributes a
\textbf{Shared Event} to all tables: a piece of news or consequence of
another table's actions that enters each group's narrative.

\begin{enumerate}[nosep]
  \item \textbf{The Rumour} (after 45 min) --- news that changes the picture
  \item \textbf{The Crisis} (after 90 min) --- something goes wrong at city scale
  \item \textbf{The Bells} (after 150 min) --- the Santo Spirito procession
    is about to begin
\end{enumerate}

\subsection{The Telegram Bot}

\nota[Bot Commands]{
  \texttt{/score [faction] [0-4]} --- update faction score\\
  \texttt{/event [text]} --- send shared event to all tables\\
  \texttt{/total} --- show aggregate score (coordinator only)\\
  \texttt{/status} --- show current status of all tables
}

\secsep

\section{Optional Physical Components}

\textit{Vespri 1282} plays perfectly well with paper and pencil. But for
a live session --- especially a multi-table evening event --- a few
physical components make the mechanics tangible. They are entirely
optional and change no rules: they only make visible what is already
happening.

\nota[Suggested Props]{
  \textbf{Stress tokens} --- nine dark stones per player, placed on the
  Stress boxes of the sheet. The sheet stays the reference; the tokens
  make the weight a character carries visible.\\[4pt]
  \textbf{Passion tokens} --- five red beads per player, for those using
  the Passion layer (see ``The Passion of the Vespers''). Spending Fortune,
  resisting, or entering a Crisis become gestures the whole table sees.\\[4pt]
  \textbf{Connection ribbons} --- one ribbon or cord for each declared
  Passion Connection. A ribbon stretched between two tables signals to
  the whole room that something binds them.
}

Two token colours are enough: Fortune spends Stress, so no third type is
needed. Keep the ribbon scissors at the coordinator's station, not at the
tables: severing a bond should be a deliberate moment.

\filettodoppio
